\documentclass[12pt]{article}
\usepackage{sectsty,setspace}
\usepackage[top=1.87cm, bottom=1.87cm, left=1.87cm, right=1.87cm]{geometry} 
\usepackage{epstopdf}
\usepackage{amsmath,latexsym,amssymb,wasysym}
\usepackage{times}
\usepackage{fancyhdr}
\usepackage[T1]{fontenc}
\usepackage[skip=10pt plus1pt, indent=0pt]{parskip} % ligne vide entre paragraphes = espace automatique, plus besoin de \\

\setstretch{1}
%\setlength{\baselineskip}{0.167in}
\setlength\parindent{0pt}

\pagestyle{fancy}
\fancyhf{}
\fancyhead[R]{Christophe Rouleau-Desrochers}
\fancyfoot[C]{\thepage}
\fancyfoot[L]{Présentation intégrée du parcours}
\renewcommand{\headrulewidth}{0pt}

\begin{document}
\section*{Présentation intégrée du parcours}
J'ai grandi sur une ferme laitière où j'ai aidé mon père dès mon plus jeune âge. Lorsqu'il est devenu acériculteur pendant mon adolescence, j'ai passé beaucoup de temps à travailler en forêt et à m'en émerveiller. C'est à cette époque qu'est né mon intérêt pour les sciences naturelles. Dans une érablière, tout est affaire de calendrier : la coulée dépend des cycles de gel et de dégel, le débourrement annonce la fin du temps des sucres et les couleurs d'automne indiquent qu'il faut rentrer le bois de chauffage. Cette temporalité si harmonieusement orchestrée par la forêt, que je ne savais pas encore nommer, allait devenir le cœur de ma recherche : la phénologie.

L'école a longtemps été un défi pour moi, et c'est d'abord par la pratique que j'ai appris. Dès l'âge de 15 ans, j'ai travaillé dans des restaurants, dont cinq à six ans en cuisine, tout en suivant une formation de trois ans en gestion de la restauration à l'ITHQ~; j'ai ensuite été sommelier pendant environ trois ans. Ce milieu exigeant m'a inculqué la rigueur et le sens du travail d'équipe, qui m'ont servi plus tard à diriger une équipe de terrain et à former 15 étudiantes et étudiants à UBC. En 2019, je suis parti travailler six mois dans un vignoble en Nouvelle-Zélande. Ce travail avec les plantes m'a convaincu de changer de voie, et en 2020, j'ai quitté la restauration pour me consacrer à l'écologie.

Je suis entré à l'université plus tard que la plupart des étudiants. J'ai d'abord terminé un certificat en écologie à l'UQAM, qui m'a permis de reprendre le chemin des études, puis un baccalauréat en biologie, où je me suis spécialisé en écologie. Je l'ai obtenu avec une moyenne de 4,30 sur 4,30 et la Médaille académique d'argent du Gouverneur général. Ornithologue amateur, j'ai mené deux projets de recherche sur les pics dans le laboratoire de Pierre Drapeau (UQAM), dont l'un, financé par une bourse de recherche de premier cycle (BRPC) du CRSNG, portait sur 20 ans de phénologie de nidification du Grand Pic. Ces projets ont été ma première immersion dans la recherche en écologie forestière.

Désirant revenir aux arbres, j'ai ensuite contacté de ma propre initiative la Dre Elizabeth Wolkovich, de l'Université de la Colombie-Britannique (UBC), pour faire un stage dans son laboratoire, financé par une deuxième BRPC. C'est là que la phénologie des arbres et l'impact des changements climatiques sur les communautés forestières sont devenus mon principal intérêt de recherche, et que j'ai cosigné avec elle un article sur la modélisation du terroir viticole. J'y ai réalisé ma maîtrise, récemment terminée avec distinction. Mon mémoire relie la phénologie foliaire et la croissance radiale de quatre espèces de bétulacées en jardin commun, issues de quatre provenances dont Saint-Hippolyte, et de 11 espèces d'arbres matures que j'ai carottés à l'Arboretum Arnold de l'Université Harvard. Des saisons plus longues ont augmenté la croissance des jeunes arbres, mais pas celle des arbres matures. Pour ces analyses, j'ai appris la modélisation bayésienne hiérarchique dans Stan, jusqu'à construire un modèle conjoint allométrie--croissance. J'ai aussi mené à terme une expérience de 630 arbres sur la longueur de la saison de croissance, dont l'analyse formera le premier chapitre de ma thèse.

Mon projet de doctorat vise à expliquer pourquoi l'allongement de la saison de croissance ne se traduit pas toujours par une croissance accrue des arbres. Il me ramènera aussi dans les forêts du Québec, à Saint-Hippolyte, où, à l'aide de drones et de dendromètres, je relierai phénologie et croissance à l'échelle de l'arbre dans une forêt mixte. C'est encore le calendrier de l'érablière que j'étudie, dans un climat plus chaud et plus incertain.

\end{document}
